\documentclass[10pt,a4paper]{amsart}
\usepackage{amsmath,amssymb,amsthm,graphicx}
% Retained minimal preamble for the Stacks Project chapter sites-cohomology.tex (Cohomology on Sites).
% Verbatim excerpt of the chapter preamble: Xy-pic setup (native preamble lines 14-18) and the native
% theorem declarations with the mathematical macros (native preamble lines 154-217). Copied mechanically;
% nothing added, removed or redefined. The wrapper class and the amsmath/amssymb/amsthm/graphicx base
% packages are supplied by the assembler, so the class/externaldocument part of the native preamble is omitted.
% For commutative diagrams we use Xy-pic
\usepackage[all]{xy}

% We use 2cell for 2-commutative diagrams.
\xyoption{2cell}
% Theorem environments.
%
\theoremstyle{plain}
\newtheorem{theorem}[subsection]{Theorem}
\newtheorem{proposition}[subsection]{Proposition}
\newtheorem{lemma}[subsection]{Lemma}

\theoremstyle{definition}
\newtheorem{definition}[subsection]{Definition}
\newtheorem{example}[subsection]{Example}
\newtheorem{exercise}[subsection]{Exercise}
\newtheorem{situation}[subsection]{Situation}

\theoremstyle{remark}
\newtheorem{remark}[subsection]{Remark}
\newtheorem{remarks}[subsection]{Remarks}

\numberwithin{equation}{subsection}

% Macros
%
\def\lim{\mathop{\mathrm{lim}}\nolimits}
\def\colim{\mathop{\mathrm{colim}}\nolimits}
\def\Spec{\mathop{\mathrm{Spec}}}
\def\Hom{\mathop{\mathrm{Hom}}\nolimits}
\def\Ext{\mathop{\mathrm{Ext}}\nolimits}
\def\SheafHom{\mathop{\mathcal{H}\!\mathit{om}}\nolimits}
\def\SheafExt{\mathop{\mathcal{E}\!\mathit{xt}}\nolimits}
\def\Sch{\mathit{Sch}}
\def\Mor{\mathop{\mathrm{Mor}}\nolimits}
\def\Ob{\mathop{\mathrm{Ob}}\nolimits}
\def\Sh{\mathop{\mathit{Sh}}\nolimits}
\def\NL{\mathop{N\!L}\nolimits}
\def\CH{\mathop{\mathrm{CH}}\nolimits}
\def\proetale{{pro\text{-}\acute{e}tale}}
\def\etale{{\acute{e}tale}}
\def\QCoh{\mathit{QCoh}}
\def\Ker{\mathop{\mathrm{Ker}}}
\def\Im{\mathop{\mathrm{Im}}}
\def\Coker{\mathop{\mathrm{Coker}}}
\def\Coim{\mathop{\mathrm{Coim}}}

% Boxtimes
%
\DeclareMathSymbol{\boxtimes}{\mathbin}{AMSa}{"02}

%
% Macros for moduli stacks/spaces
%
\def\QCohstack{\mathcal{QC}\!\mathit{oh}}
\def\Cohstack{\mathcal{C}\!\mathit{oh}}
\def\Spacesstack{\mathcal{S}\!\mathit{paces}}
\def\Quotfunctor{\mathrm{Quot}}
\def\Hilbfunctor{\mathrm{Hilb}}
\def\Curvesstack{\mathcal{C}\!\mathit{urves}}
\def\Polarizedstack{\mathcal{P}\!\mathit{olarized}}
\def\Complexesstack{\mathcal{C}\!\mathit{omplexes}}
% \Pic is the operator that assigns to X its picard group, usage \Pic(X)
% \Picardstack_{X/B} denotes the Picard stack of X over B
% \Picardfunctor_{X/B} denotes the Picard functor of X over B
\def\Pic{\mathop{\mathrm{Pic}}\nolimits}
\def\Picardstack{\mathcal{P}\!\mathit{ic}}
\def\Picardfunctor{\mathrm{Pic}}
\def\Deformationcategory{\mathcal{D}\!\mathit{ef}}
\usepackage{cleveref}
\crefname{originalreference}{original reference}{original references}
\Crefname{originalreference}{Original reference}{Original references}
\makeatletter
\newlabel{lemma-cech-vanish-collection}{{03F9}{source}{}{sourcepointer-1}{}}
\newlabel{lemma-cech-vanish-collection@cref}{{[originalreference][0][]03F9}{[1][0][]source}{}{}{}}
\makeatother
\crefname{originalreference}{original reference}{original references}
\Crefname{originalreference}{Original reference}{Original references}
\makeatletter
\newlabel{lemma-injective-module-trivial-cech}{{03FC}{source}{}{sourcepointer-2}{}}
\newlabel{lemma-injective-module-trivial-cech@cref}{{[originalreference][0][]03FC}{[1][0][]source}{}{}{}}
\makeatother
\crefname{originallemma}{original lemma}{original lemmas}
\Crefname{originallemma}{Original lemma}{Original lemmas}
\makeatletter
\newlabel{sites-lemma-when-shriek}{{00XR}{source}{}{sourcepointer-3}{}}
\newlabel{sites-lemma-when-shriek@cref}{{[originallemma][0][]00XR}{[1][0][]source}{}{}{}}
\makeatother
\title{The pullback of an injective module along a continuous cocontinuous functor has no higher cohomology on objects of the source site}
\author{The Stacks Project Authors}
\date{stacks-project:sites-cohomology.tex:local-snapshot-20261011; GNU FDL 1.2 or later}
\begin{document}
\maketitle
\noindent\textit{Editorial note.} All delivered mathematical text is the author's own, copied from the pinned Stacks Project chapter file sites-cohomology.tex; no statement or proof step was rewritten, reordered or removed.\par\medskip
\noindent\textit{Editorial note.} The counted proof is the author's complete argument: the proof is written out in full: the reduction to Cech cohomology, the identification of the two Cech complexes through cocontinuity and the fibre-product compatibility, and the final vanishing from injectivity.\par\medskip
\noindent\textit{Editorial note.} The label lemma-pullback-injective-pre-limp names the counted statement, and the proof environment immediately following it in the source is delivered verbatim as its proof.\par\medskip
\noindent\textit{Editorial note.} The counted statement is carried by the author's own lemma environment in the chapter source.\par\medskip
\noindent\textit{Editorial note.} The statement and proof are copied from the pinned chapter file; the three cited inputs are kept as external pointers with their canonical Stacks tags.\par\medskip
\noindent\textit{Allowed standard prerequisites (editorial).}\begin{itemize}
\item Cech-to-derived comparison for a covering. Same chapter, the lemma deducing the vanishing of cohomology from the vanishing of all higher Cech cohomology groups of coverings. Supplies the reduction from cohomology to Cech cohomology.
\item The pullback along a morphism of topoi as the presheaf pullback. Sites, the lemma identifying g\textasciicircum{}\{-1\} with the presheaf pullback u\textasciicircum{}p for a continuous and cocontinuous functor. Supplies the identification of the Cech complexes used in the comparison.
\item Injective modules have no higher Cech cohomology. Same chapter, the lemma stating that the higher Cech cohomology of an injective module vanishes. Supplies the vanishing of the Cech groups on the image covering.
\end{itemize}\medskip
\ifdefined\hypertarget\hypertarget{sourcepointer-1}{}\fi
\noindent\textit{Source reference.} Original reference 03F9: sources/stacks-sites-cohomology/pinned-original-source.tex, source lines 1565--1565. Prerequisite pointer invoked at proof line(s) 9289; it names a standard fact of the same chapter that the counted argument uses, and every step of the counted proof built on it is retained verbatim in the delivered body.\par\medskip
\ifdefined\hypertarget\hypertarget{sourcepointer-2}{}\fi
\noindent\textit{Source reference.} Original reference 03FC: sources/stacks-sites-cohomology/pinned-original-source.tex, source lines 1816--1816. Prerequisite pointer invoked at proof line(s) 9306; it names a standard fact of the same chapter that the counted argument uses, and every step of the counted proof built on it is retained verbatim in the delivered body.\par\medskip
\ifdefined\hypertarget\hypertarget{sourcepointer-3}{}\fi
\noindent\textit{Source reference.} Original lemma 00XR: sources/stacks-chow/tags.txt, source lines 1229--1229. Prerequisite pointer invoked at proof line(s) 9303; it names a standard fact of a sibling chapter that the counted argument uses, and every step of the counted proof built on it is retained verbatim in the delivered body.\par\medskip
\begin{lemma}
\label{lemma-pullback-injective-pre-limp}
Let $u : \mathcal{C} \to \mathcal{D}$ be a continuous and cocontinuous
functor of sites. Let $g : \Sh(\mathcal{C}) \to \Sh(\mathcal{D})$
be the corresponding morphism of topoi. Let $\mathcal{O}_\mathcal{D}$
be a sheaf of rings and let $\mathcal{I}$ be an injective
$\mathcal{O}_\mathcal{D}$-module. Then
$H^p(U, g^{-1}\mathcal{I}) = 0$ for all $p > 0$ and $U \in \Ob(\mathcal{C})$.
\end{lemma}

\begin{proof}
The vanishing of the lemma follows from
Lemma \ref{lemma-cech-vanish-collection}
if we can prove vanishing of all higher
{\v C}ech cohomology groups
$\check H^p(\mathcal{U}, g^{-1}\mathcal{I})$
for any covering $\mathcal{U} = \{U_i \to U\}$ of $\mathcal{C}$.
Since $u$ is continuous, $u(\mathcal{U}) = \{u(U_i) \to u(U)\}$
is a covering of $\mathcal{D}$, and
$u(U_{i_0} \times_U \ldots \times_U U_{i_n}) =
u(U_{i_0}) \times_{u(U)} \ldots \times_{u(U)} u(U_{i_n})$.
Thus we have
$$
\check H^p(\mathcal{U}, g^{-1}\mathcal{I}) =
\check H^p(u(\mathcal{U}), \mathcal{I})
$$
because $g^{-1} = u^p$ by Sites, Lemma \ref{sites-lemma-when-shriek}.
Since $\mathcal{I}$ is an injective
$\mathcal{O}_\mathcal{D}$-module these {\v C}ech cohomology groups vanish, see
Lemma \ref{lemma-injective-module-trivial-cech}.
\end{proof}

\end{document}
